\documentclass[11pt]{article}

\usepackage{amsmath,amssymb,amsthm}
\usepackage[margin=1in]{geometry}
\usepackage{mathtools}
\usepackage{stmaryrd}   % \llbracket \rrbracket
\usepackage{booktabs}
\usepackage[hidelinks]{hyperref}
\usepackage[htt]{hyphenat}   % allow long \texttt{} file paths to break

%--- theorem environments -------------------------------------------------
\theoremstyle{plain}
\newtheorem{theorem}{Theorem}
\newtheorem{proposition}{Proposition}
\newtheorem{lemma}{Lemma}
\newtheorem{corollary}{Corollary}
\newtheorem{conjecture}{Conjecture}
\theoremstyle{definition}
\newtheorem{definition}{Definition}
\newtheorem{question}{Question}
\newtheorem{example}{Example}
\theoremstyle{remark}
\newtheorem{remark}{Remark}

%--- macros ---------------------------------------------------------------
\newcommand{\C}{\mathcal{C}}
\newcommand{\Set}{\mathbf{Set}}
\newcommand{\Vect}{\mathbf{Vec}}
\newcommand{\Vecfd}{\mathbf{Vec}_{\mathrm{fd}}}
\newcommand{\Fam}{\mathrm{Fam}}
\newcommand{\Poly}{\mathbf{Poly}}
\newcommand{\Cont}{\mathbf{Cont}}
\newcommand{\Alg}{\mathbf{Alg}}
\newcommand{\Add}{\mathrm{Add}}
\newcommand{\ext}[1]{\llbracket #1 \rrbracket}
\newcommand{\ihom}[2]{[#1,#2]}
\renewcommand{\lhd}{\vartriangleleft}
\newcommand{\op}{\mathrm{op}}
\newcommand{\pizero}{\pi_0}
\newcommand{\Hom}{\mathrm{Hom}}
\newcommand{\one}{\mathbf{1}}
\DeclareMathOperator{\Nat}{Nat}
\DeclareMathOperator{\Ext}{Ext}
\DeclareMathOperator{\End}{End}
\DeclareMathOperator{\coker}{coker}
\DeclareMathOperator{\supp}{supp}
\DeclareMathOperator{\dm}{dim}

\setlength{\parskip}{2pt}

%==========================================================================
\begin{document}

\begin{center}
{\Large\bfseries Vector-space containers ($\mathbf{VCont}$)}\\[4pt]
{\large A self-contained treatment: the biproduct collapse, its exact finite-field\\
 defect, algebraic structure, and the admissibility frontier}\\[10pt]
\begin{tabular}{rl}
\textbf{Author:} & MacBeth (Kodamai)\\
\textbf{Date:} & 2026-09-05, rev.\ 2026-09-11\\
\textbf{For review by:} & Rick\\
\textbf{Corresponds to:} & \texttt{scotmacbeth/work-in-progress} commit \texttt{c523e82}\\
\end{tabular}\\[8pt]
{\itshape\bfseries Draft --- work in progress; not yet neighbour-reviewed.
Subject matter prepared for Neil Ghani / an external requester.}
\end{center}

\noindent\rule{\textwidth}{0.4pt}
\begin{center}
\begin{minipage}{0.92\textwidth}\small
\emph{Reading guide and honesty note.} Every theorem below is tagged in prose with its
epistemic grade: \textbf{proved} (a complete pen-and-paper proof exists in the cited internal
file), \textbf{lean-verified} (machine-checked, \texttt{sorry}-free, in Lean~4), or
\textbf{OPEN} (a precisely stated conjecture that is \emph{not} settled). The headline result of
\S2 is proved and lean-verified. The frontier question of \S4 (``Question~(Q)'') is \textbf{OPEN}
and is conjectured independent of ZFC; it is presented as such and nowhere treated as settled.
This note has not yet had a neighbour review.
\end{minipage}
\end{center}
\noindent\rule{\textwidth}{0.4pt}

\medskip
\noindent
Fix a field $k$. A \emph{container} in a category $\C$ is a pair (shapes, positions): a set $S$
together with a position object $P_s\in\C$ for each $s\in S$, extending to an endofunctor
$\ext{S,P}$ of $\C$. Over $\C=\Set$ this is the classical theory of Abbott, Altenkirch and Ghani:
containers are precisely the polynomial functors, and the extension functor is a full and faithful
embedding. This note asks what survives when the positions are \emph{vector spaces} --- the base
$\C=\Vecfd$ of finite-dimensional $k$-vector spaces --- and reports the answers the programme has
established. The one-line summary: \emph{the coproduct/product distinction that makes $\Set$
containers rigid becomes a biproduct over $\Vect$, and much of what the $\Set$ theory reads off
``for free'' collapses.} Over a finite field $\mathbb F_q$ that collapse is not merely qualitative:
\S\ref{sec:defect} computes its exact size $\delta=\prod_s a_s-\prod_s(a_s-(|T|-1))$ as a closed
integer point-count.

%==========================================================================
\section{Definition of $\mathbf{VCont}$ and its extension functor}
\label{sec:def}

Write $\Vect=\Vect_k$ for all $k$-vector spaces, $\Vecfd$ for the finite-dimensional ones, and
$\Vect(P,W)$ for the hom vector space. All functors below are $\Vect$-functors: their action on
homs $\Vect(V,W)\to\Vect(FV,FW)$ is $k$-linear.

\begin{definition}[linear / vector-space container]
\label{def:vcont}
A \emph{vector-space container} (object of $\mathbf{VCont}:=\Fam(\Vect^\op)$, or of
$\Fam(\Vecfd^\op)$ when finiteness is imposed) is a pair $(S,(P_s)_{s\in S})$ with $S$ a set (the
\emph{shapes}) and each $P_s\in\Vect$ a \emph{position space}. A morphism
$(S,(P_s))\to(T,(Q_t))$ is a pair $(f,(\varphi_s))$ with $f\colon S\to T$ a function and, for each
$s\in S$, a linear map $\varphi_s\colon Q_{f(s)}\to P_s$ --- \emph{covariant on shapes, contravariant
on positions}, exactly as for $\Set$ containers. Write $n_s:=\dm P_s$.
\end{definition}

\begin{definition}[extension]
\label{def:ext}
The \emph{extension} of $(S,P)$ is the functor $\ext{S,P}\colon\Vect\to\Vect$,
\[
   \ext{S,P}\,W \;:=\; \bigoplus_{s\in S}\Vect(P_s,\,W).
\]
Choosing a basis of $P_s$ gives $\Vect(P_s,W)\cong W^{n_s}$ naturally in $W$, so
$\ext{S,P}\,W\cong\bigoplus_{s\in S} W^{\oplus n_s}$, a ``direct sum of powers''. On a morphism
$(f,(\varphi_s))$ the natural transformation is $(g_s)_s\mapsto\bigl(g_{f(s)}\circ\varphi_s\bigr)_s$.
Writing $h_P:=\Vect(P,-)$ for the corepresentable at $P$, we have
$\ext{S,P}=\bigoplus_{s\in S}h_{P_s}$ and $h_k=\Vect(k,-)\cong\mathrm{Id}$.
\end{definition}

Two features of $\Vect$ shape everything that follows.

\emph{Coproduct $=$ product $=$ biproduct.} A finite coproduct in $\Vect$ is a biproduct, hence
\emph{equals} the product. Thus although the outer sum $\bigoplus_{s\in S}$ in Definition~\ref{def:ext}
is written as a coproduct (a ``$\Sigma$ over shapes'' of the hom-spaces $\Vect(P_s,X)$), for finite
$S$ its underlying space is equally the product $\prod_{s\in S}\Vect(P_s,X)$ --- a ``$\Pi$ over
shapes''. This $\Sigma\!\rightsquigarrow\!\Pi$ coincidence is the \emph{biproduct collapse}, and it
is the single change from the set-level extension $\ext{S,P}X=\coprod_{s\in S}X^{P_s}$ (a genuine
disjoint union). It is made fully explicit in the Lean formalisation of \S\ref{sec:lean}, where the
linear extension is literally $\mathtt{LExt}\,(S,\mathrm{Pos})\,X=(s:S)\to(\mathrm{Pos}\,s\to X)$.

\emph{Finite positions distribute over $\oplus$.} If $\dm P<\infty$ then
$\Vect(P,\bigoplus_t V_t)\cong\bigoplus_t\Vect(P,V_t)$ naturally; equivalently, $h_P$ preserves
coproducts iff $P$ is finite-dimensional (a finite-support argument on a basis of $P$). This is why
we usually restrict positions to $\Vecfd$; the infinite-dimensional case is the source of the
frontier in \S\ref{sec:admiss}.

\begin{example}
\label{ex:small}
Take $S=\{a,b\}$ with $P_a=k^2$, $P_b=k$. Then
$\ext{S,P}\,W=\Vect(k^2,W)\oplus\Vect(k,W)\cong W^2\oplus W\cong W^3$, so
$\ext{S,P}\cong\mathrm{Id}^3$. Numerically $\dm\ext{S,P}(k^d)=3d$ for every $d$
(verified computationally, $d=0,\dots,4$). The two integers that produced this functor --- the
shape $S$ and the partition $3=2+1$ --- have already vanished from the answer: only the total
$N=\sum_s n_s=3$ remains. That disappearance is the subject of \S\ref{sec:collapse}.
\end{example}

%==========================================================================
\section{The collapse: $\ext{-}$ is neither full nor faithful over $\Vect$, and not injective on objects}
\label{sec:collapse}

Over $\Set$ the extension functor $\ext{-}\colon\Fam(\Set^\op)\to[\Set,\Set]$ is \textbf{full and
faithful} (Abbott--Altenkirch--Ghani): the container data $(S,P)$ can be recovered from the functor,
with $S=\ext{S,P}(1)$ the value at a one-element set. Over $\Vect$ this recovery fails at its root.

\begin{remark}[the terminal slogan dies]
The terminal object of $\Vect$ is the zero object $0$, so
$\ext{S,P}(0)=\bigoplus_s\Vect(P_s,0)=0$: the shapes are invisible at the terminal object, and the
$\Set$ slogan ``$F(1)=S$'' has no analogue.
\end{remark}

\subsection{The finite collapse (headline; proved and lean-verified)}

\begin{theorem}[finite collapse; \textbf{proved}]
\label{thm:collapse}
Let $S$ be finite and every $n_s=\dm P_s$ finite, and put $N:=\sum_{s\in S}n_s=\dm\bigl(\bigoplus_s
P_s\bigr)$. Then there is a natural isomorphism $\ext{S,P}\cong\mathrm{Id}^N$. Consequently:
\begin{enumerate}
\item $\ext{S,P}\cong\ext{S',P'}$ as functors \textbf{iff} $N=N'$; a finite vector-space container
is classified up to isomorphism of its extension by the single number $N$.
\item Since $\End(\mathrm{Id})=k$ is local and $\End(\mathrm{Id}^N)=M_N(k)$ is semiperfect,
Krull--Schmidt--Azumaya makes the decomposition into indecomposables unique; the recovered invariant
is exactly the multiplicity $N$. The shape set $S$ and the partition $N=\sum_s n_s$ are \emph{not}
recoverable.
\end{enumerate}
\end{theorem}

\noindent This is proved in \texttt{proofs/2026-08-18-linear-containers-vec.md} (Theorem~2.1;
registry \texttt{linear-containers-vec.json}, node \texttt{part1-finite-collapse}, grade
\textbf{proved}), with the co-Yoneda lemma $\Nat(h_P,G)\cong G(P)$ and the additivity of
$P\mapsto h_P$ as the engine.

\begin{theorem}[the explicit witness; \textbf{proved}, \textbf{lean-verified}]
\label{thm:witness}
The two containers
\[
   p=(\{\ast\},\,k^2)\qquad\text{and}\qquad p'=(\{1,2\},\,(k,k))
\]
both have extension $\mathrm{Id}^2$, i.e.\ both present the functor $X\mapsto X\oplus X$, yet they
are \textbf{not isomorphic as containers}: their shape sets have different cardinalities
($1\neq2$). Hence $\ext{-}$ is \emph{not injective on objects}.
\end{theorem}

\noindent Contrast $\Set$: there the analogous containers $(\{\ast\},2)$ and $(\{1,2\},(1,1))$ have
\emph{distinct} extensions $X^2$ and $2\cdot X$, separated already at $X=1$. The collapse is
precisely the failure of that separation. Theorem~\ref{thm:witness} is machine-checked; see
\S\ref{sec:lean}.

\subsection{Why: the hom formula and the failure of extensivity}

The mechanism is visible one categorical level up, on hom-sets. Both computations below are proved
in \texttt{2026-08-18-linear-containers-vec.md} (Theorems~3.1--3.3; registry nodes
\texttt{part2-hom-formula}, \texttt{part2-extensivity-crux}, both \textbf{proved}).

\begin{proposition}[hom formula; \textbf{proved}]
\label{prop:hom}
For vector-space containers $(S,P)$, $(T,Q)$,
\begin{align*}
  \Nat\bigl(\ext{S,P},\ext{T,Q}\bigr)&\;\cong\;\prod_{s\in S}\ \bigoplus_{t\in T}\Vect(Q_t,P_s),\\
  \Fam(\Vect^\op)\bigl((S,P),(T,Q)\bigr)&\;\cong\;\prod_{s\in S}\ \coprod_{t\in T}\Vect(Q_t,P_s),
\end{align*}
where $\coprod$ is the coproduct in $\Set$ (disjoint union): a container morphism assigns to each
$s$ a single choice $t=f(s)$ together with a linear map in $\Vect(Q_{f(s)},P_s)$.
\end{proposition}

\begin{theorem}[neither full nor faithful; \textbf{proved}]
\label{thm:notfull}
The functor $\ext{-}$ induces on each hom the canonical comparison
\[
   \prod_{s}\ \coprod_{t}\Vect(Q_t,P_s)\;\longrightarrow\;\prod_{s}\ \bigoplus_{t}\Vect(Q_t,P_s),
\]
componentwise the map from the disjoint union $\coprod_t$ into the biproduct $\bigoplus_t$ sending a
choice $(t,\varphi)$ to the element supported at $t$. Over $\Vect$ this map is \emph{neither
surjective nor injective}, so $\ext{-}$ is \textbf{neither full nor faithful}.
\begin{enumerate}
\item \emph{Not full} as soon as some shape $s$ admits two shapes $t\neq t'$ with
$\Vect(Q_t,P_s)\neq0\neq\Vect(Q_{t'},P_s)$: any natural transformation whose $s$-component is a
genuine linear combination $\varphi_t+\varphi_{t'}$ supported on two shapes is not in the image of a
shape-function.
\item \emph{Not faithful} as soon as $S\neq\varnothing$ and $|T|\geq2$. Over a field the zero linear
map lies in every $\Vect(Q_t,P_s)$, so the two \emph{distinct} container morphisms $(u,0)$ and
$(u',0)$ --- with $u\neq u'\colon S\to T$ and all position maps zero --- are collapsed to the
\emph{same} (zero) natural transformation. The mechanism is exactly the biproduct: $\bigoplus_t$ has a
single shared zero across all summands, so it merges the $|T|$ distinct disjoint-union-tagged zeros
$(t,0)$ into that one zero element; the comparison is injective only \emph{away from the zero maps}.
Hence $\ext{-}$ is \textbf{faithful iff $|T|\leq1$ or $S=\varnothing$}. By contrast $\Set$'s coproduct
is the disjoint union, which keeps a distinct tag on each zero, so the classical
Abbott--Altenkirch--Ghani $\ext{-}$ over $\Set$ remains faithful.
\end{enumerate}
\end{theorem}

\begin{remark}[provenance of the faithfulness correction]
The failure of faithfulness was surfaced by R.\ Langer's finite-field analogue (q-container
Cor.~5.3, Ex.~5.4); over $\mathbb F_q$ it is a finite point-count, made exact in
\S\ref{sec:defect}. (Draft note.)
\end{remark}

The point is a single symbol. Over $\Set$ the same computation gives
$\Nat(\ext{S,P},\ext{T,Q})\cong\prod_s\bigoplus_t\Set(Q_t,P_s)$ where now $\bigoplus_t=\coprod_t$ is
the coproduct \emph{in $\Set$}, i.e.\ the disjoint union: container-hom equals $\Nat$, and $\ext{-}$
is full and faithful. \textbf{$\Set$ is extensive} --- its coproduct is the disjoint union --- and
that is exactly what makes $\ext{-}$ full. \textbf{$\Vect$ is not extensive} --- its coproduct is a
biproduct --- and the fullness gap
\[
   \coprod_t \ \subsetneq\ \bigoplus_t
\]
is the precise, quantified manifestation of that failure. The object-level collapse
(Theorem~\ref{thm:collapse}) and the morphism-level fullness gap (Theorem~\ref{thm:notfull}) are
\emph{one} phenomenon: the failure of extensivity under the base change $\Set\rightsquigarrow\Vect$.

\begin{remark}[what survives; \textbf{proved}]
The collapse used both finiteness of $S$ and finite-dimensionality of positions; drop either and
content returns. With infinite $S$ and all $P_s=k$, $\ext{S,P}=\mathrm{Id}^{\oplus|S|}$ is a genuine
coproduct, not a product, of copies of $\mathrm{Id}$; and for a single infinite-dimensional $P$,
$h_P=\Vect(P,-)$ is an infinite-product functor, not built from $\mathrm{Id}$ under $\oplus$, and $P$
\emph{is} recovered from the representable (Lemma~1.2/1.3 of the source file). The finite recovery is
sharp: $\ext{-}$ determines the total space $\bigoplus_s P_s$ up to isomorphism, and nothing about
its partition into shapes (Corollary~3.4, \textbf{proved}).
\end{remark}

%==========================================================================
\section{The finite-field instance: $q$-containers and the exact collapse defect}
\label{sec:defect}

Everything in \S\ref{sec:collapse} is stated over an arbitrary field $k$ and is \emph{qualitative}:
$\ext{-}$ either is or is not full, is or is not faithful. Over a \emph{finite} field
$k=\mathbb F_q$ these failures become finite point-counts, and the size of the collapse can be
written in closed form. This is the setting of R.\ Langer's \emph{$q$-container} note (personal
communication, 2026-09-11): Definition~\ref{def:vcont} verbatim, with positions
finite-dimensional $\mathbb F_q$-spaces. His \S4 morphism count and \S5 faithfulness corollary are
the first two displays below; the exact \emph{defect} (Theorem~\ref{thm:defect}) is a joint
refinement. This subsection is thus the arithmetic shadow of the whole of \S\ref{sec:collapse}:
the one structural symbol $\coprod_t\rightsquigarrow\bigoplus_t$ becomes the one integer
$-(|T|-1)$.

Fix $\mathbb F_q$ and $q$-containers $C=(S,(P_s))$, $D=(T,(Q_t))$ with $S,T$ finite; write
$d_s=\dm P_s$, $e_t=\dm Q_t$, and $\Cont_q(C,D):=\Fam(\Vecfd^\op)\bigl((S,P),(T,Q)\bigr)$ for the
(now finite) morphism set. The extension map on this hom,
$\ext{-}\colon\Cont_q(C,D)\to\Nat(\ext C,\ext D)$, has finite source and target, and
Proposition~\ref{prop:hom} specialises to two counts:
\begin{equation}
\label{eq:counts}
   |\Cont_q(C,D)|=\prod_{s\in S}\Bigl(\sum_{t\in T}q^{\,e_t d_s}\Bigr),
   \qquad
   |\Nat(\ext C,\ext D)|=q^{\,(\sum_s d_s)(\sum_t e_t)}.
\end{equation}
The first is a product of independent per-shape choices: at each $s$ pick a target $u(s)=t$ and a
map in $\Vect(Q_t,P_s)$, of which there are $q^{e_t d_s}$. The second is co-Yoneda: with
$P:=\bigoplus_s P_s$ and $Q:=\bigoplus_t Q_t$ the extension is corepresentable,
$\ext C=\Vect(P,-)$, so $\Nat(\ext C,\ext D)\cong\Vect(Q,P)$, of cardinality
$q^{\dm Q\cdot\dm P}$.

Under that co-Yoneda bijection a morphism $(u,(f_s))$ becomes the block matrix $M\in\Vect(Q,P)$,
rows indexed by $s$ and columns by $t$, whose $s$-th row carries the single block $f_s$ in column
$u(s)$ and zero elsewhere. So as $(u,f)$ ranges over $\Cont_q(C,D)$ the induced matrix is built
\emph{row by row, independently across rows}, and each row $g_s$ ranges over the set of tuples
supported in a single column:
\[
   A_s\;=\;\bigcup_{t\in T}W_t^{(s)},\qquad
   W_t^{(s)}:=\{\,g\in\textstyle\bigoplus_{t'}\Vect(Q_{t'},P_s):g_{t'}=0\ (t'\neq t)\,\}
   \;\cong\;\Vect(Q_t,P_s).
\]
The image of $\ext{-}$ is therefore the product $\prod_s A_s$, and the only nontrivial count is
$|A_s|$: a union of $|T|$ coordinate subspaces of a direct sum, which pairwise meet only in $0$.

\begin{lemma}[union of pairwise-trivially-meeting subspaces; \textbf{proved}]
\label{lem:union}
If $W_1,\dots,W_n$ ($n\geq1$) are finite subspaces of an $\mathbb F_q$-space with
$W_i\cap W_j=\{0\}$ for $i\neq j$, then
$\bigl|\bigcup_i W_i\bigr|=\bigl(\sum_i|W_i|\bigr)-(n-1)$.
\end{lemma}

\noindent Each $W_i$ contains $0$; for $i\neq j$ the only shared vector is $0$, so the punctured
sets $W_i\setminus\{0\}$ are pairwise disjoint and
$\bigl|\bigcup_i W_i\bigr|=1+\sum_i(|W_i|-1)=(\sum_i|W_i|)-(n-1)$. Applied to the $|T|$ coordinate
subspaces $W_t^{(s)}$, with $|W_t^{(s)}|=q^{e_t d_s}$, this gives (for $|T|\geq1$)
$|A_s|=\bigl(\sum_t q^{e_t d_s}\bigr)-(|T|-1)$.

\begin{theorem}[collapse-defect formula; \textbf{proved}]
\label{thm:defect}
Write $a_s:=\sum_{t\in T}q^{\,e_t d_s}$. For $q$-containers $C=(S,(P_s))$, $D=(T,(Q_t))$ with
$|T|\geq1$,
\[
   \bigl|\,\mathrm{image}\,\ext{-}\,\bigr|=\prod_{s\in S}\bigl(a_s-(|T|-1)\bigr),
   \qquad
   \delta(C,D):=|\Cont_q(C,D)|-\bigl|\,\mathrm{image}\,\ext{-}\,\bigr|
      =\prod_s a_s-\prod_s\bigl(a_s-(|T|-1)\bigr).
\]
\end{theorem}

\noindent Combine $|\mathrm{image}\,\ext{-}|=\prod_s|A_s|$ with Lemma~\ref{lem:union} and
subtract \eqref{eq:counts}. Proved in
\texttt{proofs/2026-09-11-collapse-defect-formula.md} (Theorem~5; registry
\texttt{linear-containers-vec.json}, node \texttt{collapse-defect-formula}, grade \textbf{proved}),
brute-force verified on 15 cases over $\mathbb F_2,\mathbb F_3,\mathbb F_4$ (dims $\leq2$,
$|S|,|T|\leq3$; \texttt{scratch/collapse\_defect\_verify.py}), all exact.

The formula makes the two qualitative statements of \S\ref{sec:collapse} into arithmetic:

\begin{corollary}[faithful/full, quantified; \textbf{proved}]
\label{cor:defect}
For $|T|\geq1$: $\ext{-}$ is faithful $\iff\delta(C,D)=0\iff|T|=1$ or $S=\varnothing$ (the threshold
of Theorem~\ref{thm:notfull}, now read off as ``every factor $a_s-(|T|-1)$ equals $a_s$''); and
$\ext{-}$ is full $\iff$ every factor $A_s$ fills its space $\iff$ for each shape $s$ with $P_s\neq0$
at most one target $t$ has $Q_t\neq0$.
\end{corollary}

The single term $-(|T|-1)$ is the whole story. It counts exactly the $|T|$ distinct
disjoint-union tags $(t,0)$ that the biproduct $\bigoplus_t$ merges into its one shared zero: the
$|T|$ morphisms $(u,f{=}0)$ with different shape-functions $u\colon S\to T$ all induce the same
zero row. Over $\Set$ the corresponding union is a genuine disjoint union $\coprod_t$, which keeps
a distinct tag on every zero; the correction term vanishes, $\delta=0$, and the
Abbott--Altenkirch--Ghani representation is full and faithful. Thus the finite-field defect is the
failure of extensivity (\S\ref{sec:collapse}) rendered as an integer.

\begin{remark}[the $q\to1$ limit does \emph{not} recover $\Set$-faithfulness; \textbf{proved}]
It is tempting to read $\delta$ through a ``field with one element''. Formally $a_s\to|T|$ and
$a_s-(|T|-1)\to1$, so $|\mathrm{image}|\to1$ while $|\Cont|\to|T|^{|S|}$: the $q\to1$ degeneration
is \emph{maximal} collapse, $\delta\to|T|^{|S|}-1\neq0$, not the faithful $\Set$ case. The reason is
structural: the $\Set$ extension $V\mapsto\coprod_s\Set(P_s,V)$ is a \emph{different functor}
(a coproduct of representables), full and faithful by AAG; the linear extension replaces that
$\Set$-coproduct by the biproduct. The $q\to1$ limit reflects only that every hom-space collapses
over $\mathbb F_1$, not any set-level faithfulness. The two functors must not be conflated.
\end{remark}

The witness that surfaced the whole faithfulness correction is the smallest nontrivial instance:
$C=(\{\ast\},k)$, $D=(\{1,2\},(k,k))$ over $\mathbb F_2$ gives $|\Cont|=4$, $|\mathrm{image}|=3$,
$\delta=1$ --- two of the four morphisms (the two zero maps with $u=1$ vs.\ $u=2$) collapse to one
natural transformation (Langer, Ex.~5.4). A sample of the verified table:

\begin{center}\small
\renewcommand{\arraystretch}{1.15}
\begin{tabular}{@{}c c c c c c c@{}}
\toprule
$q$ & $C$ & $D$ & $|\Cont|$ & $|\mathrm{image}|$ & $|\Nat|$ & $\delta$\\
\midrule
$2$ & $(1)$   & $(1,1)$   & $4$  & $3$  & $4$  & $1$\\
$2$ & $(1,1)$ & $(1,1)$   & $16$ & $9$  & $16$ & $7$\\
$2$ & $(1)$   & $(1,1,1)$ & $6$  & $4$  & $8$  & $2$\\
$2$ & $(2)$   & $(1,1)$   & $8$  & $7$  & $16$ & $1$\\
$3$ & $(1,1)$ & $(1,1)$   & $36$ & $25$ & $81$ & $11$\\
$2$ & $(1,1)$ & $(2)$     & $16$ & $16$ & $16$ & $0$ \ ($|T|{=}1$)\\
\bottomrule
\end{tabular}
\end{center}

\noindent Here $C,D$ are written as the tuples of position dimensions $(d_s)$, $(e_t)$. The
$|T|=1$ row realises Corollary~\ref{cor:defect}'s faithful case ($\delta=0$); the $(2)\to(1,1)$
row is the Theorem~\ref{thm:witness} witness direction. The finite-dimensional theory over a
finite field is thus not merely qualitatively but \emph{numerically} pinned down.

%==========================================================================
\section{Algebraic structure}
\label{sec:algebra}

Composition of vector-space containers is available on the finite-dimensional locus, and it too is
shaped by the collapse.

\begin{proposition}[composition; \textbf{proved}]
\label{prop:comp}
For containers with finite-dimensional positions,
\[
   \ext{S,P}\circ\ext{T,Q}\;\cong\;\ext{\,S\times T,\ (P_s\otimes Q_t)_{(s,t)}\,},
\]
so on $\Fam(\Vecfd^\op)$ the composition product is
$(S,P)\lhd(T,Q)=(S\times T,(P_s\otimes Q_t))$, with unit $I=(\{\ast\},k)$ (extension $\mathrm{Id}$).
\end{proposition}

\noindent This is \emph{not} the $\Set$ container composition, whose shapes form the dependent sum
$\coprod_{s\in S}T^{P_s}$: linearity flattens the dependency (a linear map $P_s\to\bigoplus_t V_t$ is
a \emph{sum of components}, not a \emph{choice of branch}) into the plain product $S\times T$.
Proved in \texttt{2026-08-18-linear-containers-vec.md}, Prop.~4.1 (registry node
\texttt{part3-composition}, \textbf{proved}).

The comonoids for this product identify a clean algebraic species.

\begin{theorem}[$\lhd$-comonoids are families of algebras; \textbf{proved}]
\label{thm:comonoid}
Let $C=(S,P)\in\Fam(\Vecfd^\op)$ with $S$ an arbitrary set and each $P_s$ finite-dimensional. The
$\lhd$-comonoid structures $(\delta,\varepsilon)$ on $C$ in $(\Fam(\Vecfd^\op),\lhd,I)$ are in
natural bijection with families $(A_s)_{s\in S}$ of unital associative $k$-algebra structures
$A_s=(P_s,\mu_s,\eta_s)$, one on each position space. The bijection is
$\delta_{\mathrm{shape}}=\Delta$ (the diagonal), $\delta^\sharp_s=\mu_s$,
$\varepsilon^\sharp_s=\eta_s$. On morphisms this upgrades to an isomorphism of categories
\[
   \mathrm{Comon}_\lhd\bigl(\Fam(\Vecfd^\op)\bigr)\ \cong\ \Fam(\Alg_k^\op),
\]
the cocommutative comonoids forming the full subcategory $\Fam(\mathbf{CAlg}_k^\op)$.
\end{theorem}

\noindent Proved in \texttt{proofs/2026-08-19-vec-comonoids-algebras.md} (registry node
\texttt{part3-comonoid-algebras}, grade \textbf{proved}; hypotheses sharpened to
finite-dimensional positions only, $S$ arbitrary). The counit forces the shape comultiplication to
the diagonal (the unique $(\Set,\times)$-comonoid on $S$); contravariance turns $\delta^\sharp_s$
into a multiplication $P_s\otimes P_s\to P_s$ and $\varepsilon^\sharp_s$ into a unit $k\to P_s$;
the comonoid laws become associativity and the two-sided unit law.

\begin{remark}[the algebroid guess, refuted in finite dimension; \textbf{proved}]
Over $\Set$ the analogous fact is $\lhd$-comonoid $\cong$ small category (Ahman--Uustalu; $\Cont\cong
\Fam(\Set^\op)$). One might hope $\lhd$-comonoid over $\Vect$ $\cong$ $k$-linear category (algebroid,
Mitchell). This is \emph{false} in finite dimension: because the counit forces
$\delta_{\mathrm{shape}}=\Delta$, the comultiplication never reaches an off-diagonal block
$P_a\otimes P_b$ ($a\neq b$), so there is no hom-between-distinct-objects composition. The structure
is a disjoint family of one-object $k$-linear categories $=$ $k$-algebras. A genuine algebroid
requires double-indexed positions and a matrix (bimodule) composition
$(P\circ Q)_{a,c}=\bigoplus_b P_{a,b}\otimes Q_{b,c}$ --- the $\bigoplus_b$ restoring the extensivity
coproduct whose absence caused the strict collapse. That identification is classical
(B\'enabou; Mitchell) and is treated separately; the finite single-index result stands as above.
\end{remark}

\subsection*{Which monoidal / closed structures survive (high level)}

Three companion results from the wider ``containers over a base'' programme locate the collapse
inside the closed-structure theory. They are stated here at high level and \emph{not} re-proved; each
is proved in its own source file.
\begin{itemize}
\item \textbf{T1 (unit connectedness; proved).} For a closed symmetric monoidal cocomplete base
$\C$, the extension functor $\ext{-}\colon\Fam(\C^\op)\to[\C,\C]$ is full and faithful \emph{iff} the
monoidal unit is \emph{connected} ($\C(I,-)$ preserves coproducts). Over $\Vect$ the unit is
\emph{disconnected} (the zero object makes $|{-}|\colon\Vecfd\to\Set$ fail to preserve coproducts),
which is exactly why $\ext{-}$ is not full --- a base-general explanation of \S\ref{sec:collapse}.
\item \textbf{T2 (Dirichlet / Day tensor; proved).} The pointwise Dirichlet tensor on
$\Fam(\C^\op)$ is Day convolution, and its closedness is controlled by a familial-representability
condition that fails over both $\Vect$ and $\Vecfd$.
\item \textbf{T4-left (left closure; proved).} The composition product $\lhd$ is left-closed on the
copower-tiny (collapse) locus, where $p\lhd q=p\otimes q$; this is the locus on which
Proposition~\ref{prop:comp} lives.
\end{itemize}
(These are internal results of MacBeth's programme; grades are as recorded in memory and are cited,
not re-derived, here.)

%==========================================================================
\section{Composition and admissibility}
\label{sec:admiss}

Proposition~\ref{prop:comp} defines $\lhd$ on \emph{finite-dimensional} positions. The general
question is whether $\Fam(\Vect^\op)$ carries a composition product at all.

\begin{definition}[$\lhd$-admissibility]
$\Fam(\C^\op)$ is \emph{$\lhd$-admissible} if the image of $\ext{-}$ in $[\C,\C]$ is closed under
composition: for all $p,q$ there is $r$ with $\ext{r}\cong\ext{p}\circ\ext{q}$.
\end{definition}

\begin{proposition}[admissibility is per-object; \textbf{proved}]
\label{prop:absorptive}
Call $P\in\C$ \emph{absorptive} if $X\mapsto\ihom{P}{\ext{q}X}$ is an extension for every $q$. Then
$\Fam(\C^\op)$ is $\lhd$-admissible \textbf{iff} every object of $\C$ is absorptive. There are two
sources of absorptivity: \emph{copower-tiny} (flexible) objects, where $\ihom{P}{\coprod_t
Y_t}=\coprod_t\ihom{P}{Y_t}$; and \emph{distributive/extensive} (rigid) objects, via the $\Set$-style
mechanism.
\end{proposition}

\noindent Proposition~6.1 of \texttt{proofs/2026-09-01-gap1-setxvec-proved.md} (\textbf{proved}). In
$\Vecfd$ every object is copower-tiny (finite-dimensional $=$ finitely generated $=$ copower-tiny;
Lemma~5.1 there, \textbf{proved}), so $\Fam(\Vecfd^\op)$ is admissible on the collapse locus with
$\lhd=\otimes_k$ --- this is Proposition~\ref{prop:comp} again.

\subsection{A proved positive result: $\Set\times\Vecfd$ is admissible}

\begin{theorem}[Gap~1 inhabited; \textbf{proved}]
\label{thm:gap1}
On the finite-vec-support locus $\mathcal P\subseteq\Fam((\Set\times\Vecfd)^\op)$ (shape $S$ and
set-positions $A_s$ arbitrary; $\{s:V_s\neq0\}$ finite), the base $\C=\Set\times\Vecfd$ is
$\lhd$-admissible: for all $p,q\in\mathcal P$ there is $r=p\lhd q\in\mathcal P$ and an
\emph{explicit natural isomorphism} $\Theta\colon\ext{p\lhd q}\Rightarrow\ext{p}\circ\ext{q}$ ---
both components built as maps (the $\Set$-component a bijection with explicit inverse, the
$\Vect$-component a natural composite of hom-out-of-coproduct, tensor--hom, and finite-support
isomorphisms), \emph{not} matched merely by cardinality.
\end{theorem}

\noindent Theorem~3.1 of \texttt{proofs/2026-09-01-gap1-setxvec-proved.md} (\textbf{proved}; the map
was built and its bijectivity, invertibility, and naturality confirmed computationally with zero
failures). The base $\Set\times\Vecfd$ is non-collapse, non-cartesian, and has a disconnected unit,
so it exhibits admissibility \emph{without} unit-connectedness: it factors a rigid (extensive)
$\Set$ part against a flexible (copower-tiny) $\Vect$ part. One honesty caveat, itself proved: the
product $\lhd$ carries \emph{no canonical} monoidal structure --- the construction must choose a
concentration slot --- and the obstruction to canonicity is exactly the non-injectivity of
$\ext{-}$ on objects (Theorem~\ref{thm:witness}). A non-canonical monoidal structure does exist.

\subsection{The open frontier: is full $\Vect$ admissible? (Question (Q), \textbf{OPEN})}

For \emph{finite-dimensional} positions everything works. For \emph{full} $\Vect$ the issue is that
an infinite-dimensional position $P$ is not copower-tiny, so $\ihom{P}{\bigoplus_t(-)}$ need not be
an extension. By Proposition~\ref{prop:absorptive} the whole question reduces to a single test
object. The sharpest one is $E:=k^{(\mathbb N)}$ (countable-dimensional, free, non-tiny) against
$q=(\mathbb N,(k))$, giving $\ext{q}X=\bigoplus_{\mathbb N}X$ and

\begin{question}[the frontier; \textbf{OPEN}]
\label{q:Q}
Is
\[
   F(X)\ :=\ \Hom\bigl(E,\ \textstyle\bigoplus_{\mathbb N}X\bigr)\ =\ \prod_{\mathbb N}\Bigl(\bigoplus_{\mathbb N}X\Bigr)
\]
\[
   (\text{row-finite }\mathbb N\times\mathbb N\text{ matrices over }X)
\]
naturally isomorphic to a coproduct of representables $\bigoplus_{j}\Hom(N_j,-)$ in the additive
functor category $\Add(\Vect,\Vect)$? Equivalently: is $F$ a \emph{projective} object of
$\Add(\Vect,\Vect)$?
\end{question}

\noindent The stakes: \textbf{(Q) YES} would give an irreducible admissible base beyond
Theorem~\ref{thm:gap1}; \textbf{(Q) NO} would confirm the ``absorptive dichotomy'' conjecture (every
nice admissible base decomposes as rigid-extensive $\times$ flexible-additive). \emph{This is the one
deep open problem in the theory and we state it as unsettled.}

What \emph{is} proved is a precise homological reduction, all in ZFC, in
\texttt{proofs/2026-09-02-vec-admissibility-rigidity.md} and
\texttt{proofs/2026-09-04-Q-homological-reduction-and-target-dimension-dichotomy.md}:
\begin{itemize}
\item \textbf{Lemma R (Yoneda rigidity; proved).} Every $\alpha\in\Nat(F,\mathrm{id})$ has
$\alpha\circ\mathrm{in}_n=0$ for all but finitely many rows $n$. The engine is pure Yoneda:
$\prod_{\mathbb N}=h_E$ and $\Nat(h_E,\mathrm{id})=E=k^{(\mathbb N)}$ is finite-support. This
resurrects Baer--Specker-type rigidity \emph{over a field}, where module slenderness fails.
\item \textbf{Theorem E (homological reduction; proved, ZFC).} With $F_{\mathrm{fin}}:=\bigoplus_n
A\subseteq F$ ($A:=\bigoplus_m\mathrm{id}$) and $Q:=F/F_{\mathrm{fin}}$, for every $M$ there is an
exact sequence $0\to\coker(\rho_M)\to\Ext^1(Q,M)\to\Ext^1(F,M)\to0$. Hence
\[
   F\text{ is projective}\ \Longleftrightarrow\ \Ext^1(F,M)=0\ \text{for all }M.
\]
So (Q) is exactly the single homological question ``is $\Ext^1(F,-)\equiv0$?''.
\item \textbf{Theorem H (cardinality is neutral; proved, ZFC).} The $\mathrm{id}$-dual
$D(h_N)=\Nat(h_N,\mathrm{id})=N$ has \emph{no} $2^{\dm}$ blow-up (unlike
$\Hom_{\mathbb Z}(\mathbb Z^{(\kappa)},\mathbb Z)$), so no cardinal invariant $\dm F(X)$ or
$\dm\Nat(F,M)$ separates $F$ from a free functor. Both classical decision routes --- single-target
rigidity (all such $\Nat$-groups vanish on $Q$) and cardinality --- are thereby \emph{provably
closed}; any resolution must engage $\Ext^1(F,-)$ at an infinite-dimensional target.
\end{itemize}

\begin{remark}[status of (Q); \textbf{OPEN}, conjectured independent of ZFC]
Question~\ref{q:Q} is \textbf{not settled}. Because both the cardinality engine (Theorem~H) and the
single-target engine are provably unavailable over a field, the residual is a Whitehead-type
extension problem on the syzygies of $F$. The current \emph{conjecture} --- theorem-supported but
\emph{not} proved --- is that (Q) is \textbf{independent of ZFC}: plausibly NO under
$2^{\aleph_0}<2^{\aleph_1}$ and possibly YES under a failure of that hypothesis plus a
uniformisation. Settling (Q) requires genuine set-theoretic input (Shelah; Eklof--Mekler transported
to functor categories) and is the flagged frontier of the programme.
\end{remark}

%==========================================================================
\section{Formalisation}
\label{sec:lean}

The object-level collapse (Theorem~\ref{thm:witness}) is machine-checked in Lean~4 core, with
\emph{no} Mathlib dependency, in \texttt{lean/Containers/Containers/VecCollapse.lean}. Positions are
carried by their finite basis types ($k^2$ by \texttt{Bool}, $k$ by \texttt{Unit}), and the linear
extension is modelled as the $\Pi$-over-shapes
$\mathtt{LExt}\,(S,\mathrm{Pos})\,X=(s:S)\to(\mathrm{Pos}\,s\to X)$ --- the single
$\Sigma\!\rightsquigarrow\!\Pi$ shift that \emph{is} the biproduct collapse. The two containers
$p=(\{\ast\},k^2)$ and $p'=(\{1,2\},k)$ appear as \texttt{VecCollapse.p} and \texttt{VecCollapse.p'}.
The theorem \texttt{Containers.VecCollapse.collapseIso} exhibits a full natural isomorphism
\texttt{LinNatIso p p'} (both directions, both triangle identities, and naturality in $X$), while
\texttt{Containers.VecCollapse.shapes\_not\_equiv} certifies (\texttt{sorry}-free, and axiom-free ---
\texttt{\#print axioms} reports \texttt{[]}) that the shape sets \texttt{Unit} and \texttt{Bool}
admit no bijection. Together they are the machine-checked witness that $\ext{-}$ presents the same
functor $X\mapsto X\oplus X$ by two non-isomorphic containers --- Theorem~D(ii), \textbf{lean-verified}.
The isomorphism \texttt{collapseIso} depends only on \texttt{Quot.sound} (function extensionality).

%==========================================================================
\section{Summary and outlook}
\label{sec:summary}

\begin{center}
\small
\renewcommand{\arraystretch}{1.25}
\begin{tabular}{@{}p{0.45\textwidth} p{0.29\textwidth} p{0.18\textwidth}@{}}
\toprule
\textbf{Statement} & \textbf{Source} & \textbf{Grade}\\
\midrule
Extension $\ext{S,P}W=\bigoplus_s\Vect(P_s,W)$; $\Sigma\!\to\!\Pi$ biproduct form
  & 2026-08-18, Def.\ 0.2 & definition\\
Finite collapse $\ext{S,P}\cong\mathrm{Id}^N$; classified by $N$ alone (Thm~\ref{thm:collapse})
  & 2026-08-18, Thm~2.1 & \textbf{proved}\\
Witness $(\{\ast\},k^2)\not\cong(\{1,2\},k)$, same functor $X\mapsto X\oplus X$ (Thm~\ref{thm:witness})
  & 2026-08-18 / VecCollapse.lean & \textbf{proved}, \textbf{lean-verified}\\
$\ext{-}$ neither full nor faithful; faithful iff $|T|\leq1$ or $S=\varnothing$ (Thm~\ref{thm:notfull})
  & 2026-08-18, Thm~3.3 & \textbf{proved}\\
$\mathbb F_q$ defect $\delta=\prod_s a_s-\prod_s(a_s-(|T|-1))$, $a_s=\sum_t q^{e_t d_s}$ (Thm~\ref{thm:defect})
  & 2026-09-11, Thm~5 & \textbf{proved}\\
Composition $(S,P)\lhd(T,Q)=(S\times T,(P_s\otimes Q_t))$ (Prop~\ref{prop:comp})
  & 2026-08-18, Prop~4.1 & \textbf{proved}\\
$\lhd$-comonoids $\cong$ families of $k$-algebras (Thm~\ref{thm:comonoid})
  & 2026-08-19, Thm~7 & \textbf{proved}\\
Admissibility $\Leftrightarrow$ every object absorptive (Prop~\ref{prop:absorptive})
  & 2026-09-01, Prop~6.1 & \textbf{proved}\\
$\Set\times\Vecfd$ is $\lhd$-admissible (Thm~\ref{thm:gap1})
  & 2026-09-01, Thm~3.1 & \textbf{proved}\\
Homological reduction: $F$ proj.\ $\Leftrightarrow\Ext^1(F,-)\equiv0$ (Thm~E)
  & 2026-09-04, Thm~E & \textbf{proved} (ZFC)\\
No cardinal invariant separates $F$ from free (Thm~H)
  & 2026-09-04, Thm~H & \textbf{proved} (ZFC)\\
\textbf{(Q)} Is full $\Vect$ admissible / $F$ projective? (Question~\ref{q:Q})
  & 2026-09-02/04 & \textbf{OPEN} (conj.\ indep.\ ZFC)\\
\bottomrule
\end{tabular}
\end{center}

\medskip
\noindent\textbf{Outlook.} The theory of vector-space containers is complete and clean on the
finite-dimensional locus: the extension functor is neither full nor faithful (faithful only when the
target has $\leq1$ shape or the source is empty) and is not injective on objects, the failure is
exactly the non-extensivity of $\Vect$, comonoids are families of algebras,
and $\Set\times\Vecfd$ realises admissibility off the connected-unit pole. The single deep frontier
is Question~(Q) --- whether \emph{infinite-dimensional} linear resources can still absorb external
shape --- now sharpened to a precise, ZFC-reduced homological question ($\Ext^1(F,-)\equiv0$?) with
both classical decision routes provably closed and independence of ZFC conjectured. Settling it is
the natural next step and would either produce a genuinely new irreducible composition base or
complete the absorptive-dichotomy classification.

\vfill
\noindent\rule{\textwidth}{0.4pt}\par
\begin{sloppypar}\raggedright
{\small\emph{Draft --- work in progress. Prepared for Neil Ghani / external requester, 2026-09-05.
Internal sources: \texttt{proofs/2026-08-18-linear-containers-vec.md},
\texttt{proofs/2026-08-19-vec-comonoids-algebras.md},
\texttt{proofs/2026-09-11-collapse-defect-formula.md},
\texttt{proofs/2026-09-01-gap1-setxvec-proved.md},
\texttt{proofs/2026-09-02-vec-admissibility-rigidity.md},
\texttt{proofs/2026-09-04-Q-homological-reduction-and-target-dimension-dichotomy.md},
\texttt{lean/\allowbreak Containers/\allowbreak Containers/\allowbreak VecCollapse.lean}.
Grades per the corresponding registry files.}}
\end{sloppypar}

\end{document}
